\documentclass{article}
\usepackage[UTF8]{ctex}
\usepackage{amsmath,amssymb,amsfonts}
\usepackage{enumitem}
\usepackage{geometry}
\geometry{a4paper, margin=1in}
% === 错题本排版优化（不参与题目标识） ===
\setlength{\parindent}{0pt}
\setlength{\parskip}{0.35em plus 0.10em minus 0.08em}
\setlength{\emergencystretch}{3em}
\setlist{itemsep=0.25em, topsep=0.40em, parsep=0pt}
\allowdisplaybreaks[2]
% === 错题本排版优化结束 ===
\begin{document}
\section*{中国石油大学(北京)2025年数学分析真题}
\section*{30．中国石油大学（北京）2025年数学分析试卷}

\subsection{题目1}
1．设 \(x_{n}=\cos \frac{x}{2} \cos \frac{x}{2^{2}} \cdots \cos \frac{x}{2^{n}}\) ，求 \(\lim _{x \rightarrow 0} \lim _{n \rightarrow \infty} x_{n}\) ．

\subsection{题目2}
2 ．求极限 \(\lim _{n \rightarrow \infty} \frac{\sqrt[n]{n!}}{n}\) ．

\subsection{题目3}
3．已知 \(f(x)=\left\{\begin{array}{ll}x^{2} \sin \frac{\pi}{x}, & x>0 ; \\ A, & x=0 ; f(x) \\ a x^{2}+b, & x<0 .\end{array}\right.\) 在 \(x=0\) 处可导，求 \(A, a, b\) ．

\subsection{题目4}
4．\(a_{1}, a_{2}, \cdots, a_{n}\) 为正数，\(n>2, f(x)=\left(\frac{a_{1}^{x}+a_{2}^{x}+\cdots+a_{n}^{x}}{n}\right)^{\frac{1}{x}}\) ，求 \(\lim _{x \rightarrow 0^{+}} f(x)\) 。

\subsection{题目5}
5．求不定积分 \(\int \frac{1}{x} \sqrt{\frac{x+2}{x-2}} \mathrm{~d} x\) ．

\subsection{题目6}
6．求 \(I=\int_{C}\left(x^{2}+2 x y\right) \mathrm{d} y, C\) 表示逆时针方向上半椭圆 \(\frac{x^{2}}{a^{2}}+\frac{y^{2}}{b^{2}}=1\) ．

二．（15 分）计算 \(f(x)=\arcsin x\) 的麦克劳林公式到 \(x^{5}\) 项．

三．（15 分）求 \(f(x, y)=x^{2}-x y+y^{2}-2 x+y\) 在全平面上的最大与最小值．

四．（15 分）作极坐标变换，将二重积分 \(\iint_{D} f\left(\sqrt{x^{2}+y^{2}}\right) \mathrm{d} x \mathrm{~d} y\) 为定积分，\(D=\{(x, y) \mid 0 \leq y \leq x \leq\) \(1\}\) 。

五．（15分）叙述数列 \(\left\{x_{n}\right\}\) 收敛的柯西收敛原理，并利用它证明：若数列 \(\left\{b_{n}\right\}\) 有界，则

\[
a_{n}=\frac{b_{1}}{1 \cdot 2}+\frac{b_{2}}{2 \cdot 3}+\cdots+\frac{b_{n}}{n(n+1)}\left(n \in \mathbb{N}_{-}^{+}\right)
\]

收敛。

六．（15 分）设 \(a_{n}>0\) 且单调递减，证明：\(\sum_{n=1}^{\infty} a_{n}\) 与 \(\sum_{n=1}^{\infty} 2^{n} a_{2^{n}}\) 同敛散．

七．（15 分）讨论积分 \(I(p)=\int_{0}^{1} x^{p-1}(\ln x)^{2} \mathrm{~d} x\) 在 \(\left[p_{0},+\infty\right)\) 上的一致收敛性，\(p_{0}>0\) ．
\end{document}
